% ═══════════════════════════════════════════════════════════════
% LEHRERHINWEISE LH-08b: El tiempo (Das Wetter)
% ═══════════════════════════════════════════════════════════════
% Fach: Spanisch | Klasse: 7 | Sequenz: 08 (Vacaciones)
% Fachfarbe: #c41e3a (Karminrot)
%
% Enthält:
% - Kurzplan (Stundenverlauf)
% - Differenzierungshinweise [B]/[S]/[E]
% - Häufige Fehler + Reaktionen
% - Material-Checkliste
% - Lösungen
% ═══════════════════════════════════════════════════════════════

\documentclass[11pt, a4paper]{article}

% ═══════════════════════════════════════════════════════════════
% SW-MODUS TOGGLE
% ═══════════════════════════════════════════════════════════════
\newif\ifswmodus
\swmodusfalse    % Standard: Farbe

% ═══════════════════════════════════════════════════════════════
% PAKETE
% ═══════════════════════════════════════════════════════════════

\usepackage[utf8]{inputenc}
\usepackage[T1]{fontenc}
\usepackage[ngerman]{babel}
\usepackage{helvet}
\renewcommand{\familydefault}{\sfdefault}

\usepackage[margin=2cm]{geometry}
\usepackage{enumitem}
\usepackage{tabularx}
\usepackage{booktabs}
\usepackage{longtable}

\usepackage{xcolor}
\usepackage{amssymb}
\usepackage{tcolorbox}
\tcbuselibrary{skins}

\usepackage{fancyhdr}
\usepackage{lastpage}
\usepackage{needspace}

% ═══════════════════════════════════════════════════════════════
% FARBEN
% ═══════════════════════════════════════════════════════════════

\definecolor{spanisch-color}{HTML}{c41e3a}
\definecolor{grau-color}{HTML}{888888}
\definecolor{hellgrau-color}{HTML}{f5f5f5}
\definecolor{basis-color}{HTML}{2a7a4b}
\definecolor{standard-color}{HTML}{2c5aa0}
\definecolor{erweiterung-color}{HTML}{7b1fa2}
\definecolor{loesung-color}{HTML}{c62828}

\definecolor{sw-dark}{HTML}{000000}
\definecolor{sw-gray}{HTML}{666666}
\definecolor{sw-white}{HTML}{ffffff}

\ifswmodus
    \colorlet{spanisch}{sw-dark}
    \colorlet{grau}{sw-gray}
    \colorlet{hellgrau}{sw-white}
    \colorlet{basis}{sw-dark}
    \colorlet{standard}{sw-dark}
    \colorlet{erweiterung}{sw-dark}
    \colorlet{loesung}{sw-dark}
\else
    \colorlet{spanisch}{spanisch-color}
    \colorlet{grau}{grau-color}
    \colorlet{hellgrau}{hellgrau-color}
    \colorlet{basis}{basis-color}
    \colorlet{standard}{standard-color}
    \colorlet{erweiterung}{erweiterung-color}
    \colorlet{loesung}{loesung-color}
\fi

\newcommand{\fachfarbe}{spanisch}

% ═══════════════════════════════════════════════════════════════
% VARIABLEN
% ═══════════════════════════════════════════════════════════════

\newcommand{\thema}{El tiempo -- Das Wetter}
\newcommand{\sequenz}{08}
\newcommand{\block}{b}
\newcommand{\fach}{Spanisch}
\newcommand{\klasse}{7}
\newcommand{\dauer}{90 Min. (Doppelstunde)}
\newcommand{\lerngruppengroesse}{24}
\newcommand{\datum}{\today}

% ═══════════════════════════════════════════════════════════════
% KOPF- UND FUSSZEILE
% ═══════════════════════════════════════════════════════════════

\pagestyle{fancy}
\fancyhf{}
\renewcommand{\headrulewidth}{0.4pt}
\renewcommand{\footrulewidth}{0.4pt}

\lhead{\fach{} -- Klasse \klasse}
\chead{\textbf{LH-\sequenz\block{} (Lehrerhinweise)}}
\rhead{\datum}

\lfoot{}
\cfoot{}
\rfoot{Seite \thepage\ von \pageref{LastPage}}

% ═══════════════════════════════════════════════════════════════
% BEFEHLE
% ═══════════════════════════════════════════════════════════════

\newcommand{\B}{\textcolor{basis}{\textbf{[B]}}}
\newcommand{\Std}{\textcolor{standard}{\textbf{[S]}}}
\newcommand{\E}{\textcolor{erweiterung}{\textbf{[E]}}}

\newcommand{\loesung}[1]{\textcolor{loesung}{\textbf{#1}}}

\newcommand{\es}[1]{\textcolor{\fachfarbe}{\textbf{#1}}}

\newtcolorbox{kurzplan}{
    colback=hellgrau,
    colframe=\fachfarbe,
    colbacktitle=white,
    coltitle=\fachfarbe,
    boxrule=1pt,
    arc=0pt,
    fonttitle=\bfseries,
    attach boxed title to top left={yshift=-2mm, xshift=5mm},
    boxed title style={boxrule=0pt, colback=white},
    title=Kurzplan
}

\newtcolorbox{differenzierung}{
    colback=white,
    colframe=grau,
    colbacktitle=white,
    coltitle=grau,
    boxrule=0.5pt,
    arc=0pt,
    fonttitle=\bfseries,
    attach boxed title to top left={yshift=-2mm, xshift=5mm},
    boxed title style={boxrule=0pt, colback=white},
    title=Differenzierung
}

\newtcolorbox{fehlerbox}{
    colback=white,
    colframe=loesung,
    colbacktitle=white,
    coltitle=loesung,
    boxrule=0.5pt,
    arc=0pt,
    fonttitle=\bfseries,
    attach boxed title to top left={yshift=-2mm, xshift=5mm},
    boxed title style={boxrule=0pt, colback=white},
    title=Häufige Fehler
}

\newtcolorbox{materialbox}{
    colback=white,
    colframe=\fachfarbe,
    colbacktitle=white,
    coltitle=\fachfarbe,
    boxrule=0.5pt,
    arc=0pt,
    fonttitle=\bfseries,
    attach boxed title to top left={yshift=-2mm, xshift=5mm},
    boxed title style={boxrule=0pt, colback=white},
    title=Material-Checkliste
}

% ═══════════════════════════════════════════════════════════════
% DOKUMENT
% ═══════════════════════════════════════════════════════════════

\begin{document}

% ═══════════════════════════════════════════════════════════════
% KOPFBEREICH
% ═══════════════════════════════════════════════════════════════

\begin{center}
    {\Large\bfseries\textcolor{\fachfarbe}{Lehrerhinweise: \thema}}\\[0.3em]
    {\normalsize LH-\sequenz\block{} | \dauer}
\end{center}

\vspace{10pt}

% ═══════════════════════════════════════════════════════════════
% 1. KURZPLAN
% ═══════════════════════════════════════════════════════════════

\section*{1. Stundenverlauf (Kurzplan)}

\textbf{1. Stunde (45 Min.)}

\begin{tabularx}{\textwidth}{|c|l|X|l|}
\hline
\textbf{Zeit} & \textbf{Phase} & \textbf{Inhalt} & \textbf{Material} \\
\hline
0--5 & Einstieg & Wetterkarte Spanien zeigen, Vermutungen sammeln & Projektor \\
\hline
5--15 & Input & Wettervokabeln mit Bildern einführen, vorsprechen & Bildkarten \\
\hline
15--25 & Übung 1 & Chorisches Sprechen, Aussprachekorrektur & -- \\
\hline
25--35 & Übung 2 & Memory-Spiel (Bild-Wort) in 4er-Gruppen & Karten \\
\hline
35--45 & Sicherung & AB Merkkasten 1 ausfüllen & AB \\
\hline
\end{tabularx}

\vspace{1em}

\textbf{2. Stunde (45 Min.)}

\begin{tabularx}{\textwidth}{|c|l|X|l|}
\hline
\textbf{Zeit} & \textbf{Phase} & \textbf{Inhalt} & \textbf{Material} \\
\hline
0--5 & Wiederholung & Wetter-Bingo (L ruft Begriffe, SuS markieren) & Bingo-Karten \\
\hline
5--15 & Input & Jahreszeiten einführen, Wetter zuordnen & Tafel \\
\hline
15--25 & Übung 3 & Lernpfad Schritt 3-4 oder AB Aufg. 1-2 & Tablets/AB \\
\hline
25--35 & Anwendung & Mini-Wetterbericht schreiben (PA) & AB Aufg. 4 \\
\hline
35--45 & Abschluss & Präsentationen (2-3 Paare), HA & -- \\
\hline
\end{tabularx}

% ═══════════════════════════════════════════════════════════════
% 2. DIFFERENZIERUNG
% ═══════════════════════════════════════════════════════════════

\section*{2. Differenzierung}

\begin{differenzierung}
\textbf{Legende:} \B{} Basis (BR) | \Std{} Standard (MR) | \E{} Erweiterung (GY)

\vspace{0.5em}

\textbf{WICHTIG:} Diese Marker sind NUR für die Lehrkraft. Auf Schülermaterial erscheinen KEINE Niveaumarkierungen!
\end{differenzierung}

\vspace{1em}

\begin{tabularx}{\textwidth}{|l|X|X|X|}
\hline
& \B{} \textbf{Basis} & \Std{} \textbf{Standard} & \E{} \textbf{Erweiterung} \\
\hline
\textbf{Aufgabe 1} & Zuordnung mit Bildunterstützung & Zuordnung selbstständig & + mündliche Sätze bilden \\
\hline
\textbf{Aufgabe 2} & Lückentext mit Wortliste & Lückentext selbstständig & + eigene Beispiele \\
\hline
\textbf{Aufgabe 3} & Satzanfänge vorgeben & Stichwörter als Hilfe & Freie Satzbildung \\
\hline
\textbf{Aufgabe 4} & 3 Sätze mit Vorlage & 4-5 Sätze mit Hilfen & 6+ Sätze frei \\
\hline
\end{tabularx}

\vspace{1em}

\textbf{Konkrete Hinweise:}
\begin{itemize}
    \item \B{} SuS mit LRS: Zeitzugabe (+10 Min.), vergrößerte AB-Version, Audioversion der Vokabeln
    \item \B{} DaZ-SuS: Wetterbilder mit deutschen UND spanischen Begriffen, Visualisierung verstärken
    \item \E{} Leistungsstarke: Moderatorenrolle bei Präsentationen, zusätzlichen Wetterbericht für anderen Ort
\end{itemize}

% ═══════════════════════════════════════════════════════════════
% 3. HÄUFIGE FEHLER
% ═══════════════════════════════════════════════════════════════

\section*{3. Häufige Fehler}

\begin{fehlerbox}
\begin{tabularx}{\textwidth}{|l|X|X|}
\hline
\textbf{Fehler} & \textbf{Beispiel} & \textbf{Reaktion} \\
\hline
Aussprache \textit{llueve} & *[ljueve] statt \es{['llwebe]} & Zungenposition zeigen, ll = Zunge an Gaumen \\
\hline
``Es ist kalt'' wörtlich & *\textit{Es frío} statt \es{Hace frío} & Merkspruch: ``Das Wetter MACHT kalt'' \\
\hline
Verwechslung hace/está & *\textit{Hace nublado} statt \es{Está nublado} & Regel: hace + Nomen, está + Adjektiv \\
\hline
Genus bei Jahreszeiten & *\textit{el primavera} statt \es{la primavera} & Eselsbrücke: Frühling ist die Dame \\
\hline
Akzent vergessen & *\textit{otono} statt \es{otoño} & Auf Tildetaste hinweisen \\
\hline
\end{tabularx}
\end{fehlerbox}

\vspace{1em}

\textbf{Typische Interferenzen Deutsch $\rightarrow$ Spanisch:}
\begin{itemize}
    \item Deutsche Satzstellung: *\textit{Hoy hace en Barcelona sol} $\rightarrow$ \es{Hoy hace sol en Barcelona}
    \item ``Es'' als Subjekt: Im Spanischen entfällt das Subjektpronomen bei unpersönlichen Verben
\end{itemize}

% ═══════════════════════════════════════════════════════════════
% 4. MATERIAL-CHECKLISTE
% ═══════════════════════════════════════════════════════════════

\section*{4. Material-Checkliste}

\begin{materialbox}
\begin{itemize}[label=$\square$]
    \item AB-08b.pdf (\lerngruppengroesse{} Kopien)
    \item ML-08b.pdf (1× für Lehrkraft)
    \item Lehrwerk: Qué pasa 1, S. 108--110
    \item Bildkarten Wetter (16 Stück: je 8 Bild + 8 Wort)
    \item Memory-Karten (6 Sets für 4er-Gruppen)
    \item Bingo-Karten Wetter (Klassensatz, 24 Stück)
    \item Wetterkarte Spanien (Projektion/Ausdruck A3)
    \item Tafelmarker / Kreide
    \item Optional: LP-08b.html (Tablets/PCs, mind. 12 Geräte für PA)
\end{itemize}
\end{materialbox}

% ═══════════════════════════════════════════════════════════════
% 5. LÖSUNGEN
% ═══════════════════════════════════════════════════════════════

\section*{5. Lösungen AB-08b}

\textbf{Merkkasten 1: El tiempo}

\begin{tabular}{ll}
\es{Hace sol.} & $\rightarrow$ \loesung{Die Sonne scheint.} \\
\es{Hace calor.} & $\rightarrow$ \loesung{Es ist heiß/warm.} \\
\es{Hace frío.} & $\rightarrow$ \loesung{Es ist kalt.} \\
\es{Hace viento.} & $\rightarrow$ \loesung{Es ist windig.} \\
\es{Hace buen tiempo.} & $\rightarrow$ \loesung{Das Wetter ist gut.} \\
\es{Hace mal tiempo.} & $\rightarrow$ \loesung{Das Wetter ist schlecht.} \\
\es{Llueve.} & $\rightarrow$ \loesung{Es regnet.} \\
\es{Nieva.} & $\rightarrow$ \loesung{Es schneit.} \\
\es{Está nublado.} & $\rightarrow$ \loesung{Es ist bewölkt.} \\
\end{tabular}

\vspace{1em}

\textbf{Merkkasten 2: Las estaciones}

\begin{tabular}{ll}
\es{la primavera} & $\rightarrow$ \loesung{der Frühling} \\
\es{el verano} & $\rightarrow$ \loesung{der Sommer} \\
\es{el otoño} & $\rightarrow$ \loesung{der Herbst} \\
\es{el invierno} & $\rightarrow$ \loesung{der Winter} \\
\es{¿Qué tiempo hace?} & $\rightarrow$ \loesung{Wie ist das Wetter?} \\
\end{tabular}

\vspace{1em}

\textbf{Aufgabe 1: Zuordnung} (4 Punkte)

\begin{tabular}{ll}
\es{Hace sol.} & $\rightarrow$ \loesung{Die Sonne scheint.} \\
\es{Hace frío.} & $\rightarrow$ \loesung{Es ist kalt.} \\
\es{Hace viento.} & $\rightarrow$ \loesung{Es ist windig.} \\
\es{Llueve.} & $\rightarrow$ \loesung{Es regnet.} \\
\end{tabular}

\vspace{1em}

\textbf{Aufgabe 2: Lückentext} (4 Punkte)

\begin{tabular}{ll}
a) & \loesung{calor} \\
b) & \loesung{nieva} \\
c) & \loesung{nublado} \\
d) & \loesung{primavera} \\
\end{tabular}

\vspace{1em}

\textbf{Aufgabe 3: Jahreszeiten} (4 Punkte)

Mögliche Antworten:
\begin{itemize}
    \item a) \loesung{En invierno hace frío y nieva. / En invierno hace mal tiempo.}
    \item b) \loesung{En otoño hace viento y llueve. / En otoño está nublado.}
    \item c) \loesung{En primavera hace buen tiempo y hace sol.}
\end{itemize}

\textit{Bewertung: Je 1P für korrekten Wetterausdruck + sinnvoller Kontext}

\vspace{1em}

\textbf{Aufgabe 4: Mini-Wetterbericht} (8 Punkte)

Bewertungskriterien:
\begin{itemize}
    \item Ort genannt (Estoy en...): 1P
    \item Jahreszeit genannt (Es verano...): 1P
    \item Mindestens 2 Wetterausdrücke: 2P
    \item Aktivität/Vorhaben (Voy a...): 2P
    \item Sprachliche Korrektheit: 2P
\end{itemize}

Beispiellösung: \loesung{¡Hola! Estoy en Barcelona. Es verano y hace mucho calor. Hoy hace sol y hace buen tiempo. Voy a la playa porque hace calor.}

% ═══════════════════════════════════════════════════════════════
% 6. HAUSAUFGABE
% ═══════════════════════════════════════════════════════════════

\section*{6. Hausaufgabe}

\begin{itemize}
    \item Lehrwerk S. 109, Aufgabe 3 (Wetterdialog ergänzen)
    \item Vokabeln lernen: hace sol, hace calor, hace frío, hace viento, llueve, nieva, está nublado, la primavera, el verano, el otoño, el invierno
    \item Optional: Lernpfad LP-08b.html vollständig bearbeiten (digital)
\end{itemize}

% ═══════════════════════════════════════════════════════════════
% 7. REFLEXIONSNOTIZEN
% ═══════════════════════════════════════════════════════════════

\section*{7. Reflexionsnotizen (nach Durchführung)}

\begin{tabularx}{\textwidth}{|l|X|}
\hline
\textbf{Was lief gut?} & \\[2em]
\hline
\textbf{Was lief nicht?} & \\[2em]
\hline
\textbf{Zeitmanagement} & \\[2em]
\hline
\textbf{Für nächstes Mal} & \\[2em]
\hline
\end{tabularx}

\end{document}
